% 2011 年北京卷理科第 4 题：程序框图。
% 来源 PDF：普通高考/2011/2011北京理.pdf，SHA-256 79351b2b5d5afd29a2d46d50526284443edad4deb64fbaf841db31ce2570c311。
% 原页：第 1 页（A3 三栏）第 1 栏下部，书页码 647。
% 重绘依据：照原页竖排主干 开始→i=0,s=2→判断 i<4；「否」向下→输出 s→结束；
%   「是」向右→i=i+1→向上→s=(s-1)/(s+1)→向左回到主干（箭头落在主干上），再向下进入判断框。
\begin{tikzpicture}[>=stealth, semithick,
  every node/.style={inner sep=2pt},
  box/.style={draw, minimum height=0.68cm, minimum width=2.9cm, align=center},
  round/.style={draw, rounded corners=6pt, minimum height=0.68cm,
    minimum width=1.9cm, align=center},
  para/.style={draw, trapezium, trapezium left angle=70, trapezium right angle=110,
    minimum height=0.68cm, minimum width=2.2cm, align=center},
  dec/.style={draw, diamond, aspect=1.9, minimum width=2.4cm,
    minimum height=1.02cm, align=center},
  arr/.style={->}]

  \node[round] (start) at (0,1.45) {开始};
  \node[box]   (init)  at (0,0.15) {$i=0$，$s=2$};
  \node[dec]   (test)  at (0,-3.05) {$i<4$};
  \node[para]  (out)   at (0,-4.65) {输出 $s$};
  \node[round] (end)   at (0,-5.95) {结束};

  \node[box, minimum width=3.3cm, minimum height=1.45cm]
    (sbox) at (4.05,-1.35) {$s=\dfrac{s-1}{s+1}$};
  \node[box] (inc) at (4.05,-3.05) {$i=i+1$};

  \draw[arr] (start) -- (init);
  \draw[arr] (init) -- (test);
  \draw[arr] (test) -- node[right, pos=0.3] {否} (out);
  \draw[arr] (out) -- (end);

  % 「是」分支
  \draw[arr] (test.east) -- node[above] {是} (inc.west);
  \draw[arr] (inc.north) -- (sbox.south);
  % 回边：箭头落在主干上
  \draw[arr] (sbox.west) -- (0,-1.35);
\end{tikzpicture}
